\section{初步概念}\label{sec:15}
\subsection{带奇点的纤维化}\label{sec:15a}
我们将考察映射 $\pi:E\to B$，其中 $E$、$B$ 为带角正合辛流形，$B$ 连通，$\dim(E)\geq\dim(B)$，且 $\pi$ 为 $(I_E,I_B)$ 全纯映射。另作三个假设，它们共同保证映射沿边界具有合理的行为。
\begin{itemize}
\item \emph{与 $\partial B$ 横截。}令 $C\subset\partial B$ 为任意边界层。对每个满足 $y=\pi(x)\in C$ 的 $x\in E$，有 $D\pi(TE_x)+TC_y=TB_y$。
\end{itemize}
这显然蕴含 $\pi^{-1}(C)$ 为 $E$ 的边界层，余维与 $C$ 相同。所有这样的层之并为 $\partial^vE=\pi^{-1}(\partial B)$，称为 $E$ 的\emph{垂直边界}。所有不包含在 $\partial^vE$ 中的边界面之并，称为 $E$ 的\emph{水平边界}，记为 $\partial^hE$。
\begin{itemize}
\item \emph{沿 $\partial^hE$ 正则。}若 $F$ 是 $E$ 的一个不包含在 $\partial^vE$ 中的边界面，则 $\pi|F:F\to B$ 为光滑纤维化。
\end{itemize}
这蕴含每个纤维 $E_y$ 在其边界 $\partial E_y=E_y\cap\partial^hE$ 附近光滑。此前集中讨论微分拓扑，但还须处理辛方面的性质。在任意点 $x\in E$，切空间分解为
\begin{equation}\label{eq:15.1}
TE_x=TE_x^h\oplus TE_x^v,
\end{equation}
其中垂直部分 $TE_x^v=\ker(D\pi_x)$，水平部分 $TE_x^h$ 为其关于 $\omega_E$ 的正交补；等价地，也可取关于 $\omega_E$ 与 $I_E$ 所确定度量的正交补。我们在\secref{14c}中曾短暂遇到这样的分解，但在那里它只起技术作用。最后一项要求是
\begin{itemize}
\item \emph{$\partial^hE$ 的水平性。}若 $x$ 位于某个不包含在 $\partial^vE$ 中的边界面 $F$ 上，则 $TE_x^h\subset TF_x$。
\end{itemize}
满足以上三个条件的映射 $\pi:E\to B$，称为\emph{（正合辛）带奇点的纤维化}。

考察几个简单推论。记 $\operatorname{Crit}(\pi)\subset E$ 为临界点集，$\operatorname{Critv}(\pi)=\pi(\operatorname{Crit}(\pi))\subset B$ 为临界值集。如上所述，沿 $\partial^hE$ 正则意味着 $\operatorname{Crit}(\pi)\subset E\setminus\partial^hE$。当 $\dim(B)=2$ 时，与 $\partial B$ 横截还蕴含 $\operatorname{Crit}(\pi)\cap\partial^vE=\varnothing$；但底空间维数更高时就不再如此。

\SourcePage{221}
\StatementNumber{1}\begin{lemma}\label{st:15.1}
若 $y\notin\operatorname{Critv}(\pi)$，则纤维 $E_y$，配备限制得到的 $\omega_{E_y}=\omega_E|E_y$、$\theta_{E_y}=\theta_E|E_y$、$I_{E_y}=I_E|E_y$，是带角正合辛流形。
\end{lemma}
\begin{proof}
唯一需要非平凡验证的，是负 Liouville 向量场沿 $\partial E_y=E_y\cap\partial^hE$ 指向内部。已知 $E$ 的负 Liouville 向量场沿每个边界面 $F$ 严格向内。现假设 $F$ 不属于 $\partial^vE$。由于 $TE^h$ 平行于 $F$，该向量场的 $TE^v$ 分量也沿 $F$ 严格向内。而这一分量恰是 $E_y$ 的负 Liouville 向量场，从而得到结论。
\end{proof}

分解\eref{15.1}在 $B\setminus\operatorname{Critv}(\pi)$ 上定义一个联络。由于 $TE^h$ 平行于 $\partial^hE$，该联络具有良定义的平行移动映射；即对每条路径 $\beta:[0,1]\to B\setminus\operatorname{Critv}(\pi)$，得到一个选定的微分同胚
\begin{equation}\label{eq:15.2}
h_\beta:E_{\beta(0)}\longrightarrow E_{\beta(1)},
\end{equation}
它光滑依赖于路径的选择。标准计算表明，$h_\beta$ 实际上是带角正合辛流形的同构。固定 $\pi$ 的一个正则值 $*$，记相应纤维为 $M=E_*$，令 $B'\subset B\setminus\operatorname{Critv}(\pi)$ 为 $*$ 的一个小邻域。利用适当的一族平行移动映射，可得 $B'$ 上的局部平凡化 $E'=B'\times M\to E$。拉回 $\theta_E,\omega_E$ 得到的微分形式 $\theta_{E'},\omega_{E'}$ 具有形式
\begin{equation}\label{eq:15.3}
\theta_{E'}=\theta_M+\kappa+d(\text{某个函数}),\qquad
\omega_{E'}=\omega_M+d\kappa,
\end{equation}
其中一形式 $\kappa$ 在 $(TE')^v$ 方向上消失。可将其视为底空间上取值于函数的一形式
\[
\kappa\in\Omega^1(B',C^\infty(M,\R)).
\]
事实上，从函数转为 Hamilton 向量场之后，$\kappa$ 恰是 $(TE')^h$ 的联络一形式。\znote{213}{由\eref{15.3}，扭曲完全由 $\kappa$ 记录：$\omega_{E'}$ 与 $\omega_M$ 之差恰为 $d\kappa$；而 $\kappa$ 在垂直方向上消失，故 $d\kappa$ 在每条纤维上的限制为零，它只记录不同纤维之间的比较。$\kappa$ 闭（局部为恰当形式）时，平凡化就是正合辛同构；一般情形下 $d\kappa$ 度量水平分布与辛结构的失配。这也解释了为何下文把“沿纤维恒定的变形”实现为改变联络：只须指定新的 $\kappa$，而不改变每条纤维的辛形式与原始形式。[来源：原书 §(15a) 与 §(15b)，印刷页 213–215。]}为看出这一点，取 $B'$ 上向量场 $Y$，令 $Y^h=(Y,X)$ 为其到 $E'$ 的唯一水平提升。在 Cartan 公式 $i_{Y^h}(d\kappa)=\mathcal L_{Y^h}\kappa-d(i_{Y^h}\kappa)$ 中，Lie 导数仍是纤维方向上消失的一形式。因此，$Y^h$ 的定义性质可写为
\[
i_X\omega_M-d(i_Y\kappa)=0\qquad\text{在每个纤维 }\{y\}\times M\text{ 上}.
\]
所以，$X$ 在每个纤维上都是函数 $-\kappa(Y)$ 的 Hamilton 向量场，证明了所声称的 $\kappa$ 的性质。上述考察表明，在 $B$ 的正则部分上，我们得到的是带选定联络的 $\operatorname{Aut}(M)$ 纤维丛。注意，$B$ 上的辛结构在这里不起作用。

\StatementNumber{2}\begin{note}\label{st:15.2}
文献\bcite{131}对带奇点的纤维化的定义，把我们的第三个条件，即 $\partial^hE$ 的水平性，换成下述更强的版本：
\SourcePage{222}
\begin{itemize}
\item \emph{$\partial^hE$ 附近的平凡性。}令 $M=E_*$ 为某个纤维。则存在 $B\times\partial M$ 的开邻域 $V'\subset B\times M$、$\partial^hE$ 的开邻域 $V\subset E$，以及逐纤维的微分同胚 $V'\to V$，在 $V'\cap(\{*\}\times M)$ 上为恒等，使 $\theta_E$ 的拉回为 $\theta_B+\theta_M$，$I_E$ 的拉回为乘积殆复结构 $I_B\times I_M$。
\end{itemize}
对 $\theta_E$ 的条件意味着，联络 $TE^h$ 在 $\partial^hE$ 附近平凡，并且在底空间的正则部分上给出结构群到 $\operatorname{Aut}(M,\partial M)$ 的典范约化。这一性质使纤维化稍易处理，但也有自然构造不满足它。
\end{note}

\subsection{同构与变形}\label{sec:15b}
若两个带奇点的纤维化的总空间及底空间分别存在带角正合辛流形同构，并与纤维化映射相容，则称这两个纤维化同构。同样，显然可以定义这类纤维化的一般变形概念，即所涉及全部几何数据的单参数族。这里不展开一般定义，而集中于两类较特殊的变形；稍后陈述若干技术结果时会用到它们。第一类的主要特征是保持正则纤维的同构类，第二类则完全不改变纤维的结构。

固定流形 $E,B$ 及其间映射 $\pi$。在它们上取正合辛结构的光滑族 $(\omega_E^s,\theta_E^s,I_E^s)$、$(\omega_B^s,\theta_B^s,I_B^s)$，$s\in[0,1]$，使 $\pi$ 对每个 $s$ 都成为带奇点的纤维化。若存在 $E$ 上光滑函数 $R^s$，使由
\begin{equation}\label{eq:15.4}
\omega_E^s(\cdot,Y^s)=\frac{\partial}{\partial s}\theta_E^s-dR^s
\end{equation}
定义的向量场 $Y^s$ 平行于 $\partial^hE$，则称该变形\emph{沿 $\partial^hE$ 行为良好}。等价地，取乘积 $\widetilde\pi=\Id\times\pi:\widetilde E=[0,1]\times E\to\widetilde B=[0,1]\times B$，并在总空间上赋予 $\omega_{\widetilde E}=d(\theta_E^s+R^s\,ds)$。这显然在 $\ker(D\widetilde\pi)$ 上为辛，故得到如\eref{15.1}的 $T\widetilde E$ 分解；\eref{15.4}是说，该分解的水平部分 $T\widetilde E^h$ 平行于 $\partial^h\widetilde E=[0,1]\times\partial^hE$。这允许对 $[0,1]\times(B\setminus\operatorname{Critv}(\pi))$ 中的路径定义平行移动映射。特别地，沿 $[0,s]\times\{y\}$ 移动，得到正合辛同构
\begin{equation}\label{eq:15.5}
k_y^s:(E_y,\omega_{E_y}^0,\theta_{E_y}^0)\longrightarrow
(E_y,\omega_{E_y}^s,\theta_{E_y}^s).
\end{equation}
更具体地，它们是通过积分 $Y^s|E_y$ 关于 $\omega_E^s$ 的垂直分量定义的；若在 Moser 论证中，用 $\partial_s\theta_{E_y}^s-dR^s|E_y$ 作为 $\partial_s\omega_{E_y}^s$ 的原始形式，则使用的正是这个向量场。$k_y^s$ 光滑依赖于 $(s,y)$，并在 Hamilton 同痕意义下，与沿 $B\setminus\operatorname{Critv}(\pi)$ 中路径的平行移动交换。换言之，$E$ 的正则部分的纤维丛结构，在沿 $\partial^hE$ 行为良好的变形下不变。

\SourcePage{223}
另一种限制更强的概念，是\emph{沿纤维恒定的变形}。如前，这由单参数族 $(\omega_E^s,\theta_E^s,I_E^s)$ 和 $(\omega_B^s,\theta_B^s,I_B^s)$ 给出，但现在要求一形式 $\partial_s\theta_E^s-dR^s$ 在 $TE^v$ 上消失。显然，这意味着\eref{15.5}中的映射是恒等。此外，还要求在 $\partial^hE$ 的某个邻域内，$I_E^s|TE^v$ 与 $s$ 无关。

取带奇点的纤维化 $\pi:E\to S$，其中 $S$ 为带边界曲面。令 $\kappa$ 为 $E$ 上一形式，在 $TE^v$ 方向上消失，且在 $\partial E\cup\operatorname{Crit}(\pi)$ 的某个邻域中为零。对 $s\in[0,1]$ 及某个常数 $c\gg0$，考虑
\begin{equation}\label{eq:15.6}
\begin{aligned}
\theta_S^s&=(cs+1)\theta_S,&\omega_S^s&=d\theta_S^s,& I_S^s&=I_S,\\
\theta_E^s&=\theta_E+s\kappa+cs\pi^*\theta_S,&\omega_E^s&=d\theta_E^s.
\end{aligned}
\end{equation}
由于 $\omega_E^s|TE^v=\omega_E|TE^v$ 非退化，得到依赖于 $s$ 的分解 $TE_x=TE_x^{h,s}\oplus TE_x^v$，它是\eref{15.1}的变形。取 $c$ 充分大，可保证 $\omega_E^s|TE_x^{h,s}$ 关于底空间诱导的定向为正，从而 $\omega_E^s$ 为辛。取 $I_E^s$ 为唯一的 $\omega_E^s$ 相容殆复结构，在 $TE^v$ 方向等于 $I_E$，在 $TE^{h,s}$ 上则与 $I_S$ 的拉回一致。由于 $\partial E$ 附近有 $TE^{h,s}=TE^h$，$I_E^s$ 在该处与 $I_E$ 一致，故边界弱 $I_E^s$ 凸。其次，$\theta_E^s,\omega_E^s$ 关联的负 Liouville 向量场沿 $\partial^hE$ 严格向内，因为该性质只依赖于直和项 $TE^v$；见\stref{引理}{15.1}的证明。现在考虑边界圆周 $C\subset\partial S$ 及相应的 $F=\pi^{-1}(C)\subset\partial^vE$。对 $x\in F$，$TE_x^h$ 包含一维子空间 $TF_x\cap TE_x^h$；将其投影到 $S$ 并使用该处的边界定向，就得到典范定向。负 Liouville 向量场沿 $F$ 指向内部，等价于
\begin{equation}\label{eq:15.7}
\theta_E|(TF_x\cap TE_x^h)>0.
\end{equation}
由于 $\theta_S$ 本身使 $S$ 成为带角正合辛流形，有 $\pi^*\theta_S|(TF_x\cap TE_x^h)>0$。结合 $\partial^vE$ 附近的等式 $\theta_E^s=\theta_E+cs\pi^*\theta_S$，这意味着\eref{15.7}的对应条件在整个变形中持续成立。现在已经验证，$(E,\omega_E^s,\theta_E^s,I_E^s)$ 对每个 $s$ 都是带角正合辛流形，故\eref{15.6}是沿纤维恒定的变形。从纤维丛观点看，其效果是通过 $\kappa$ 改变辛联络。为稍后使用，记下一个具体应用：

\StatementNumber{3}\begin{lemma}\label{st:15.3}
令 $\pi:E\to S$ 为带奇点的纤维化，$S$ 为带边界曲面。令 $\beta:[0,1]\to\operatorname{int}(S)\setminus\operatorname{Critv}(\pi)$ 为嵌入路径，$(\phi^s)$ 为 $\operatorname{Aut}(E_{\beta(0)},\partial E_{\beta(0)})$ 中从 $\phi^0=\id$ 开始的同痕。则存在沿纤维恒定的变形，在 $s=0$ 时为给定纤维化结构，使沿 $\beta$ 所得的平行移动映射族满足
\[
h_\beta^s=h_\beta\circ\phi^s.
\]
\SourcePage{224}
此外，若开子集 $U\subset S$ 包含 $\beta$ 的某一部分，则可安排变形，使 $\partial_s(\theta_E^s-\pi^*\theta_S^s)$ 在 $\pi^{-1}(U)$ 之外消失。
\end{lemma}
